\documentclass[usenames,dvipsnames,letterpaper,landscape,final]{slides}

\def\fileHome{.}
\def\icons{\fileHome/icons}

\input mySlideStyle2019-1.tex
\def\myfboxBlack#1{\fboxrule2pt\fbox{$\ds#1$}}
\newcommand\topSeparatorHeading[1]{\noPagetrue\begin{myslide}\headingline\vbox to5.5in{%
\vfil\begin{center}{{\large\myfboxBlack{\myfboxBlack{%
\text{\ \ \bf\Blue{\begin{tabular}{c}#1\end{tabular}\ \ }}}}}}%
\end{center}\vfil}\end{myslide}\noPagefalse}

\usepackage{multirow}
\usepackage{enumitem}
\usepackage[ruled,vlined]{algorithm2e}
%\usepackage{wallpaper}
		  \usepackage{mathtools}

\newcounter{outlineNumber}\setcounter{outlineNumber}{0}
\def\outno{\stepcounter{outlineNumber}\Blue{\theoutlineNumber. }}

\definecolor{boxColor}{rgb}{.4, .05, .2}

%------------------------------------------------------------------------------
\tikzset{cross/.style={cross out, draw=black, minimum size=2*(#1-\pgflinewidth), inner sep=0pt, outer sep=0pt},cross/.default={1pt}}

\renewcommand\floatpagefraction{0}
\renewcommand\textfraction{0}
\renewcommand\topfraction{1}
\renewcommand\bottomfraction{1}
\setcounter{topnumber}{10}
\setcounter{bottomnumber}{10}
\setcounter{totalnumber}{10}

\def\picWidth#1#2{\includegraphics[clip=true,keepaspectratio=true,width=#1]{#2}}
\def\picHeight#1#2{\includegraphics[clip=true,keepaspectratio=true,height=#1]{#2}}
\def\picWidthHeight#1#2#3{\includegraphics[clip=true,width=#1,height=#2]{#3}}

%\let\itemOld=\item
%\def\item{\parskip=0pt\itemsep=0pt\itemOld}


\begin{document}

% ===========================================

\begin{myslide}
\headingline
		  %\vspace*{-30pt}
		  {\bf\Large\begin{center}
					 Simulation of Partial Melting Mantle
		  \end{center}}

\vfill

\begin{center}{\obeylines
Chenyu Tian
\bigskip
%		  {\it\OliveGreen{CSEM} 
%		  \Blue{\it Oden Institute}}
%\BurntOrange{The University of Texas at Austin}}
	  {CSEM, Oden Institute}
The University of Texas at Austin}
\end{center}

\vfill

{\tiny\begin{center}
					 Todd Arbogast, Oden Institute \& Mathematics, University of Texas at Austin\\
Marc A. Hesse, Oden Institute \& Geosciences, University of Texas at Austin\\
		  \end{center}}

		  \vspace*{40pt}

\end{myslide}

% ===========================================

\begin{myslide}
\heading{Introduction}

\Oblue{Mantle melting}:
\begin{nlist}
\item Mantle melting is a crucial geophysical process in lithospheric evolution and plate tectonics.
\item Geophysical evidence indicates that the asthenosphere is often partially molten.
\item A thorough understanding of partial melting is fundamental to comprehending the dynamics of 
magmatism. 
\end{nlist}

\vspace{20pt}

\Oblue{Partial Melting}:
		  \begin{nlist}
		  \item Modeled as a two-phase flow framework.
		  \item Coupled with complicated chemical interactions.
		  \item Compaction is essential in deciding melt extraction and segregation.
		  \end{nlist}

\vspace{20pt}

		  {\Oorange{Challenge}}:
\begin{nlist}
\item Degenercy caused by no-melt regions poses mathematical and numerical difficulties.  
\end{nlist}

\end{myslide}

% ===========================================

\begin{myslide}
\heading{New framework}

\Oblue{Model}
\begin{nlist}
\item A degenerate two-phase Darcy–Stokes system describing the coupled interaction
between the solid matrix and the interstitial melt.
\item A nonlinear system governing the transport of chemical species and thermal
energy.
\item A eutectic phase relation characterizing the melting processes and connecting
the transport and static mechanical system through enthalpy method.
\end{nlist}

\vspace{20pt}
\Oblue{Numerics}
\begin{nlist}
\item $H(\text{div})$-and $H^1$- conforming mixed finite element spaces for Darcy and Stokes subproblems.
\item A more versatile and flexible multilevel WENO weighting scheme to address the nonlinear transport of chemical species and energy in the 
presence of porous media.
\end{nlist}

\end{myslide}

% ===========================================
\begin{myslide}
\heading{Darcy-Stokes System}

		  \Oblue{Full static system}:
\begin{alignat}{2}
    \label{eq:darcy}
		  &\phi_l (\mathbf{v}_l - \mathbf{v}_s)
		  + \frac{k_{\phi_l}}{\mu_l}(\nabla p_l - \rho_f\mathbf{g})
	    &&= 0,\\
	\label{eq:continuous}
		  &\nabla \cdot (\phi_l\mathbf{v}_l + \phi_s\mathbf{v}_s) &&=0,\\
	\label{eq:stokes}
		  &\nabla\bar p - \nabla \cdot \pmb{\tau}(\mathbf{v}_s) + \overline{\rho}\mathbf{g}&&= 0,\\
	\label{eq:compaction}
		  &\mu_s\nabla \cdot \mathbf{v}_s - \phi_l(p_l - p_s) &&=0, 
\end{alignat}

		  \vspace{20pt}
		  \Oblue{Deviatoric stress}: 
\begin{equation}\label{eq:deviatoricstress}
		  \pmb{\tau}(\mathbf{v}_s)
		  = 2\mu_s\phi_s\big(\mathcal{D} \mathbf{v}_s - \tfrac{1}{3}\nabla \cdot \mathbf{v}_s \mathbf{I}\big),
\end{equation}

\end{myslide}

\begin{myslide}
\heading{}

\end{myslide}

\end{document}
